% ============================================================
%  TIF_PLMB — Documentación del Repositorio
%  Autora: Maria Corina Hernandez
% ============================================================

\documentclass[12pt, a4paper]{article}

% --- Idioma y codificación ---
\usepackage[utf8]{inputenc}
\usepackage[T1]{fontenc}
\usepackage[spanish, es-tabla]{babel}
\usepackage{amsmath}

% --- Geometría ---
\usepackage[top=2.5cm, bottom=2.5cm, left=2.5cm, right=2.5cm]{geometry}

% --- Tipografía ---
\usepackage{lmodern}
\usepackage{microtype}

% --- Colores ---
\usepackage[dvipsnames]{xcolor}
\definecolor{uniandes}{RGB}{155, 0, 20}       % rojo institucional Uniandes
\definecolor{codebg}{RGB}{245, 245, 245}
\definecolor{notecolor}{RGB}{230, 240, 255}
\definecolor{noteborder}{RGB}{100, 130, 200}

% --- Hipervínculos ---
\usepackage[colorlinks=true, linkcolor=uniandes, urlcolor=uniandes,
            citecolor=uniandes, bookmarks=true]{hyperref}

% --- Tablas ---
\usepackage{booktabs}
\usepackage{tabularx}
\usepackage{array}
\usepackage{longtable}
\newcolumntype{L}[1]{>{\raggedright\arraybackslash}p{#1}}
\newcolumntype{C}[1]{>{\centering\arraybackslash}p{#1}}

% --- Listas ---
\usepackage{enumitem}
\setlist[itemize]{noitemsep, topsep=4pt}

% --- Código fuente ---
\usepackage{listings}
\lstset{
  backgroundcolor=\color{codebg},
  basicstyle=\ttfamily\small,
  breaklines=true,
  frame=single,
  framesep=4pt,
  rulecolor=\color{gray!40},
  xleftmargin=6pt,
  xrightmargin=6pt,
  aboveskip=8pt,
  belowskip=8pt,
  showstringspaces=false,
  inputencoding=utf8,
  extendedchars=true,
  literate={á}{{\'a}}1 {é}{{\'e}}1 {í}{{\'i}}1 {ó}{{\'o}}1 {ú}{{\'u}}1
           {ñ}{{\~n}}1 {Á}{{\'A}}1 {É}{{\'E}}1 {Í}{{\'I}}1 {Ó}{{\'O}}1
           {Ú}{{\'U}}1 {Ñ}{{\~N}}1 {ü}{{\"u}}1 {¿}{{\textquestiondown}}1
           {¡}{{\textexclamdown}}1 {–}{{--}}1 {—}{{---}}1
}

% --- Cuadros de nota / advertencia ---
\usepackage{tcolorbox}
\tcbuselibrary{skins}
\newtcolorbox{notebox}{
  colback=notecolor, colframe=noteborder,
  arc=3pt, boxrule=0.6pt,
  left=6pt, right=6pt, top=4pt, bottom=4pt,
  fonttitle=\bfseries\small,
}

% --- Árbol de directorios ---
\usepackage{dirtree}

% --- Secciones ---
\usepackage{titlesec}
\titleformat{\section}{\large\bfseries\color{uniandes}}{\thesection.}{0.5em}{}[\titlerule]
\titleformat{\subsection}{\normalsize\bfseries}{\thesubsection.}{0.5em}{}
\titleformat{\subsubsection}{\normalsize\itshape}{\thesubsubsection.}{0.5em}{}

% --- Encabezado / pie de página ---
\usepackage{fancyhdr}
\pagestyle{fancy}
\fancyhf{}
\lhead{\textcolor{gray}{\small TIF\_PLMB — Documentación del Repositorio}}
\rhead{\textcolor{gray}{\small M.C. Hernandez}}
\cfoot{\thepage}
\renewcommand{\headrulewidth}{0.3pt}

% --- Fecha ---
\usepackage{datetime2}

% ============================================================
\begin{document}
% ============================================================

% --- Portada ---
\begin{titlepage}
  \centering
  \vspace*{2cm}
  {\color{uniandes}\rule{\linewidth}{1.5pt}}\\[0.4cm]
  {\LARGE\bfseries Análisis de Factibilidad TIF para la PLMB}\\[0.3cm]
  {\large Modelación Econométrica y Geoespacial en Stata}\\[0.2cm]
  {\color{uniandes}\rule{\linewidth}{1.5pt}}\\[1.5cm]
  {\large\itshape Documentación del Repositorio \texttt{TIF\_PLMB}}\\[2cm]
  {\large Maria Corina Hernandez}\\[0.3cm]
  {\normalsize \href{mailto:mc.hernandezl12@uniandes.edu.co}{mc.hernandezl12@uniandes.edu.co}}\\[0.5cm]
  {\normalsize Universidad de los Andes — Grupo Sur}\\[2cm]
  {\normalsize Febrero–Marzo 2026}\\[1cm]
  \vfill
  {\small\textcolor{gray}{Repositorio: \href{https://github.com/mcorinahl/TIF_PLMB}{github.com/mcorinahl/TIF\_PLMB}}}
\end{titlepage}

\tableofcontents
\newpage

% ============================================================
\section{Descripción del Proyecto}
% ============================================================

Este repositorio contiene el código fuente para estimar el mayor valor del avalúo catastral en el área de influencia de la Primera Línea del Metro de Bogotá (PLMB). El objetivo central es evaluar la viabilidad del instrumento \textbf{Tax Increment Financing (TIF)} para financiar la extensión de la línea, utilizando un panel de datos catastrales de la UAECD (2014–2025).

La estrategia de identificación es un modelo de \textbf{Diferencias en Diferencias (DiD)} que compara la evolución de los avalúos catastrales entre predios dentro y fuera del área de influencia del metro, antes y después del anuncio/contratación del proyecto (2019).

% ============================================================
\section{Estructura del Repositorio}
% ============================================================

\dirtree{%
.1 TIF\_PLMB/.
.2 Código/.
.2 Datos/.
.3 raw/.
.4 INFORMACIÓN CATASTRAL Y PREDIAL/ \DTcomment{PREDIO\_FINAL\_20YY.csv, 2014–2025}.
.4 SHAPES/ \DTcomment{Shapefiles y DBFs de ArcGIS}.
.5 Lotes\_catastrales.gdb.zip.gdrive/.
.5 Lotes\_catastrales\_treatment/ \DTcomment{Buffers 400/800/1200m}.
.5 Lotes\_catastrales\_treat\_robust/.
.5 manzanas\_contrafac.csv.
.5 manzanas\_linea\_interv.csv.
.5 manzanas\_AI\_vias\_tramos.csv.
.4 EM2021\_800m\_PLMB/ \DTcomment{Datos de empleo 2021}.
.4 Nueva Área de Influencia/ \DTcomment{BASE\_NUEVA\_AI\_20YY.csv, 2014–2025}.
.3 processed/ \DTcomment{Archivos .dta generados por los do-files}.
.3 outcomes/ \DTcomment{Tablas (.docx) y gráficos (.pdf) finales}.
.2 inflacion.xlsx \DTcomment{IPC mensual 2014–2025 (DANE)}.
.2 README.md.
}

% ============================================================
\section{Estructura del Código}
% ============================================================

El flujo de trabajo se divide en nueve etapas principales.

% -----------------------------------------------------------
\subsection{Etapa 1 — Preparación de Datos (\texttt{1\_append\_merge.do})}
% -----------------------------------------------------------

Limpieza y estandarización de las 12 bases catastrales anuales de la UAECD. Crea un panel balanceado (solo predios observados los 12 años) y lo vincula con la información geográfica de tratamiento proveniente de ArcGIS.

\begin{itemize}
  \item \textbf{Entrada:} \texttt{PREDIO\_FINAL\_20YY.csv} (2014–2025), \texttt{Lotes\_catastrales.dbf}
  \item \textbf{Salida:} \texttt{predios\_completos.dta}, \texttt{predios\_balanceado.dta}, \texttt{predios\_tratamiento\_shp.dta}, \texttt{predios.dta}
\end{itemize}

% -----------------------------------------------------------
\subsection{Etapa 2 — Modelo DiD Principal (\texttt{2\_analysis\_did.do})}
% -----------------------------------------------------------

Implementación del modelo de Diferencias en Diferencias. Deflacta los avalúos a precios constantes de 2014 usando el IPC de diciembre de cada año, genera el logaritmo natural de los valores reales, y estima el efecto promedio del tratamiento (ATT) con efectos fijos de lote y año. Incluye un DiD dinámico (ventanas 2014–2025) para verificar tendencias paralelas.

\begin{itemize}
  \item \textbf{Entrada:} \texttt{predios.dta}, \texttt{inflacion.xlsx}
  \item \textbf{Salida:} \texttt{predios\_proc.dta}, \texttt{DID\_simple.docx} (parcial), \texttt{DiD\_windows\_1200.pdf}, \texttt{DiD\_windows\_com\_1200.pdf}, \texttt{DiD\_windows\_filter.pdf}
\end{itemize}

% -----------------------------------------------------------
\subsection{Etapa 3a — Atributos Físicos (\texttt{3\_analysis\_did\_atributes.do})}
% -----------------------------------------------------------

Analiza si el tratamiento está asociado a cambios en atributos físicos de los predios (área construida, número de pisos), tanto a nivel de predio como de manzana. Incluye un diseño de discontinuidad espacial con el Área de Influencia de Manzanas (AIM).

\begin{itemize}
  \item \textbf{Entrada:} \texttt{predios\_proc.dta}, \texttt{manzanas\_linea\_interv.csv}, \texttt{manzanas\_uso\_suelo.dta}, \texttt{manzanas\_treat\_aim.csv}, \texttt{manzanas\_untreat\_aim\_donut.csv}
  \item \textbf{Salida:} \texttt{manzanas\_proc.dta}, \texttt{DID\_atributes.docx}, \texttt{DID\_atributes\_aim.docx}, \texttt{DiD\_windows\_area.pdf}
\end{itemize}

% -----------------------------------------------------------
\subsection{Etapa 3b — Efectos Heterogéneos (\texttt{3\_analysis\_did\_heter.do})}
% -----------------------------------------------------------

Estima efectos heterogéneos del tratamiento según destino económico del predio (residencial, comercial, industrial, dotacional, etc.), estrato socioeconómico y tramo de la línea. Los destinos y estratos se fijan al periodo pre-tratamiento para evitar endogeneidad por cambio de uso.

\begin{itemize}
  \item \textbf{Entrada:} \texttt{predios\_proc.dta}
  \item \textbf{Salida:} \texttt{DID\_heter\_dest.docx}, \texttt{DID\_heter\_estr.docx}, \texttt{DID\_heter\_dest\_res.docx}, \texttt{DID\_heter\_tram.docx}, \texttt{DiD\_windows\_tram.pdf}
\end{itemize}

% -----------------------------------------------------------
\subsection{Etapa 4 — Coarsened Exact Matching (\texttt{4\_analysis\_cem.do})}
% -----------------------------------------------------------

Aplica CEM para mejorar la comparabilidad entre grupos de tratamiento y control, usando como covariables de \textit{matching}: estrato, área construida, número de pisos, destino económico y distancias al CBD, Transmilenio, malla vial y trazado PLMB (todas en quintiles). El \textit{matching} se realiza sobre el corte transversal de 2018 (último año pre-tratamiento). Los pesos se fusionan al panel completo para la estimación DiD ponderada.

\begin{itemize}
  \item \textbf{Entrada:} \texttt{predios\_proc.dta}, archivos \texttt{Distancias\_Manzanas-*.csv}
  \item \textbf{Salida:} \texttt{treat\_CEM.dta}, \texttt{treat\_weights.dta}, \texttt{DID\_CEM.docx}
\end{itemize}

% -----------------------------------------------------------
\subsection{Etapa 5 — Análisis Exploratorio: Distintos Buffers (\texttt{5\_robustez.do})}
% -----------------------------------------------------------

Corre el modelo DiD para buffers de 400m y 1200m usando una segunda definición de shapefiles de tratamiento (\texttt{Lotes\_catastrales\_treat\_robust}). Este ejercicio fue realizado para explorar sensibilidad a distintos radios por solicitud del equipo; no constituye la especificación base.

\begin{itemize}
  \item \textbf{Entrada:} \texttt{predios\_proc.dta}, \texttt{Lotes\_catastrales\_treat\_robust.dbf}
  \item \textbf{Salida:} \texttt{predios\_robust.dta}, \texttt{DID\_simple.docx} (columnas 400m y 1200m)
\end{itemize}

% -----------------------------------------------------------
\subsection{Etapa 6 — Nueva Área de Influencia (\texttt{6\_nueva\_ai.do})}
% -----------------------------------------------------------

Replica el pipeline completo usando la \textbf{nueva área de influencia} como grupo de tratamiento. A diferencia de un buffer de radio fijo, esta área está delimitada por avenidas y vías importantes cercanas al corredor del metro. Fue sugerida por entidades gubernamentales para dar soporte legal a la delimitación: respeta la geometría urbana (no corta manzanas a la mitad) y es la \textbf{definición por defecto para todos los análisis posteriores a este archivo}.

El grupo de control es el universo catastral de Bogotá fuera del área de influencia (ciudad completa menos el grupo de tratamiento).

\begin{itemize}
  \item \textbf{Entrada:} \texttt{BASE\_NUEVA\_AI\_20YY.csv} (2014–2025), \texttt{predios\_completos.dta}, \texttt{Lotes\_catastrales.dbf}, \texttt{inflacion.xlsx}, \texttt{manzanas\_AI\_vias\_tramos.csv}
  \item \textbf{Salida:} \texttt{predios\_proc\_ai2.dta}, columnas ``AI Vías'' en \texttt{DID\_heter\_dest.docx}, \texttt{DID\_heter\_estr.docx}, \texttt{DID\_heter\_tram.docx}
\end{itemize}

% -----------------------------------------------------------
\subsection{Etapa 7 — Índices de Accesibilidad Urbana (\texttt{7\_centralidades.do})}
% -----------------------------------------------------------

Construye índices de accesibilidad urbana a nivel de manzana. Genera tres tipos de índices:

\begin{itemize}
  \item \textbf{Accesibilidad de transporte:} promedio simple, tipo gravedad ($\gamma=0.5$) y decaimiento exponencial ($\lambda=0.01$) de distancias inversas a vías troncales, locales e intermedias.
  \item \textbf{Acceso económico:} $\log(\text{empleo}) \times \text{índice de transporte}$.
  \item \textbf{Equipamientos (amenidades):} índice de distancias inversas a equipamientos de educación, salud y recreación.
  \item \textbf{Índice compuesto $I_{AccB}$:} análisis de componentes principales (PCA) sobre 17 variables de distancia/accesibilidad, con ponderación por eigenvalor; produce 4 dimensiones (accesibilidad estructural, periferia-ruralidad, capilaridad urbana, corredores específicos).
\end{itemize}

Los índices se pegan a \texttt{predios\_proc\_ai2.dta} para su uso en análisis de heterogeneidad.

\begin{itemize}
  \item \textbf{Entrada:} \texttt{predios\_proc\_ai2.dta}, \texttt{Distancias\_Manzanas-*.csv}, \texttt{base\_bogota\_270221.dta}
  \item \textbf{Salida:} \texttt{manzanas\_access.dta}
\end{itemize}

% -----------------------------------------------------------
\subsection{Etapa 8 — Robustez: Grupo de Control Alternativo (\texttt{8\_robustez\_control.do})}
% -----------------------------------------------------------

Prueba de sensibilidad a la definición del grupo de control. En lugar de usar toda la ciudad como control (especificación base), define un contrafactual espacialmente acotado: las manzanas dentro de un \textbf{buffer de 1500m alrededor de la nueva área de influencia} pero por fuera del área de tratamiento.

\begin{itemize}
  \item \textbf{Entrada:} \texttt{predios\_proc\_ai2.dta}, \texttt{manzanas\_contrafac.csv}, \texttt{manzanas\_AI\_vias\_tramos.csv}
  \item \textbf{Salida:} Estimaciones DiD con control restringido al anillo perimetral
\end{itemize}

% -----------------------------------------------------------
\subsection{Etapa 9 — Heterogeneidad por Destino/Estrato $\times$ Accesibilidad (\texttt{9\_heter\_access.do})}
% -----------------------------------------------------------

Estima efectos heterogéneos de tercer orden: $\text{treat} \times \text{destino económico} \times \text{accesibilidad}$ y $\text{treat} \times \text{estrato} \times \text{accesibilidad}$. Responde si el efecto diferencial del metro sobre cada tipo de predio varía según el nivel de accesibilidad urbana preexistente de la manzana. Relevante para TIF: identifica si la base imponible capturable se concentra en zonas de alta o baja accesibilidad.

\begin{itemize}
  \item \textbf{Entrada:} \texttt{predios\_proc\_ai2.dta}, \texttt{manzanas\_access.dta}
  \item \textbf{Salida:} \texttt{DID\_heter\_dest\_access.docx}, \texttt{DID\_heter\_estr\_access.docx}, \texttt{DiD\_heter\_dest\_access.pdf}
\end{itemize}

% ============================================================
\section{Metodología}
% ============================================================

\subsection{Especificación Principal}

\begin{lstlisting}[language={}]
reghdfe ln_avaluo_real_2014 treat, a(codigo_lote year) vce(cluster codigo_barrio)
\end{lstlisting}

\begin{itemize}
  \item \textbf{Variable dependiente:} \texttt{ln\_avaluo\_real\_2014} — logaritmo del avalúo catastral en pesos constantes de 2014.
  \item \textbf{Variable de tratamiento:} \texttt{treat} $= 1$ si el predio está dentro del área de influencia \textbf{y} el año es $\geq 2019$. La variable concreta (\texttt{treatment\_800}, \texttt{treatment\_1200} o la variable de la nueva AI) depende del objetivo del análisis; la definición \textbf{base es 800m} (\texttt{treatment\_800}).
  \item \textbf{Efectos fijos:} lote (\texttt{codigo\_lote}) $+$ año (\texttt{year}).
  \item \textbf{Errores estándar:} clusterizados a nivel de barrio (\texttt{codigo\_barrio}).
  \item \textbf{Grupo de control:} universo de predios de Bogotá fuera del área de influencia (ciudad completa menos el grupo de tratamiento), excepto en \texttt{8\_robustez\_control.do} donde se usa un contrafactual espacialmente definido.
\end{itemize}

\subsection{DiD Dinámico (Tendencias Paralelas)}

Para verificar el supuesto de tendencias paralelas y visualizar la dinámica temporal del efecto, se estima un \textbf{DiD dinámico} (o DiD con ventanas de tiempo): se interactúa el indicador de tratamiento con dummies de año calendario, omitiendo 2018 como año de referencia.

\begin{notebox}
\textbf{Nota terminológica.} Los gráficos de coeficientes resultantes (\texttt{DiD\_windows\_*.pdf}) suelen denominarse coloquialmente \textit{event study plots} en la literatura, pero técnicamente son \textbf{coeficientes de un DiD dinámico indexados por año calendario} (2014, 2015, \ldots, 2025), no por tiempo relativo al tratamiento ($t-5, t-4, \ldots$). Dado que el tratamiento es simultáneo para todas las unidades (2019), ambas representaciones son matemáticamente equivalentes; solo diferiría el etiquetado del eje $x$. En un diseño con adopción escalonada (\textit{staggered}), la distinción sería sustantiva.
\end{notebox}

\begin{lstlisting}[language={}]
forvalues i = 2014/2025 {
    gen trat_`i' = cond(year==`i', 1*treatment_800, 0)
}
replace trat_2018 = 0   // anio omitido (referencia)

reghdfe ln_avaluo_real_2014 trat_*, a(codigo_lote year) vce(cluster codigo_barrio)
coefplot ..., keep(trat_*) vertical yline(0)
\end{lstlisting}

\subsection{Deflactación}

Los valores nominales se convierten a precios constantes de 2014 usando el IPC de diciembre de cada año (DANE):

\[
\text{cum\_mult}_{2014} = 1, \qquad
\text{cum\_mult}_{t} = \text{cum\_mult}_{t-1} \times \left(1 + \frac{\text{IPC}_t}{100}\right) \quad \forall\, t > 2014
\]
\[
\texttt{avaluo\_real\_2014} = \frac{\texttt{valor\_avaluo\_nominal}}{\text{cum\_mult}_t}
\]

\subsection{Exclusiones}

Se excluyen de todas las estimaciones los predios con los siguientes destinos económicos (no capturable por TIF):

\begin{itemize}
  \item VIAS, ESPACIO PÚBLICO, LOTE DEL ESTADO, DOTACIONAL PÚBLICO, RECREACIONAL PÚBLICO
\end{itemize}

\subsection{Áreas de Influencia}

\begin{table}[h!]
\centering
\begin{tabular}{@{}L{3.5cm} L{3.5cm} L{7cm}@{}}
\toprule
\textbf{Definición} & \textbf{Variable} & \textbf{Uso} \\
\midrule
400m & \texttt{treatment\_400} & Exploratoria (buffer estrecho) \\
\textbf{800m} & \textbf{\texttt{treatment\_800}} & \textbf{Especificación base} \\
1200m & \texttt{treatment\_1200} & Ejercicio exploratorio (solicitado; no es la especificación principal) \\
Nueva AI & \texttt{treatment} (archivo 6) & \textbf{Especificación por defecto para análisis posteriores} \\
\bottomrule
\end{tabular}
\caption{Definiciones del área de influencia utilizadas en el proyecto.}
\end{table}

\begin{notebox}
\textbf{Nueva Área de Influencia.} No es un buffer de radio fijo. Está delimitada por avenidas y vías importantes cercanas al corredor del metro, lo que la hace coincidir aproximadamente con el buffer de 800m pero respeta la geometría urbana real (manzanas y predios completos). Esta definición fue sugerida por entidades gubernamentales y le da soporte legal al perímetro de influencia, reduciendo la arbitrariedad de un radio fijo que puede cortar manzanas y predios a la mitad.
\end{notebox}

% ============================================================
\section{Datos Procesados Intermedios}
% ============================================================

\begin{longtable}{@{}L{4.5cm} L{6cm} L{4cm}@{}}
\toprule
\textbf{Archivo} & \textbf{Descripción} & \textbf{Generado en} \\
\midrule
\endhead
\texttt{predios\_completos.dta} & Panel completo sin balancear & \texttt{1\_append\_merge.do} \\
\texttt{predios\_balanceado.dta} & Solo predios observados los 12 años & \texttt{1\_append\_merge.do} \\
\texttt{predios\_tratamiento\_shp.dta} & Variables de tratamiento y tramo del shapefile & \texttt{1\_append\_merge.do} \\
\texttt{predios.dta} & Panel balanceado + tratamiento & \texttt{1\_append\_merge.do} \\
\texttt{predios\_proc.dta} & Panel + valores reales deflactados ($\sim$5.4 GB) & \texttt{2\_analysis\_did.do} \\
\texttt{predios\_robust.dta} & Panel + tres definiciones de buffer ($\sim$5.5 GB) & \texttt{5\_robustez.do} \\
\texttt{manzanas\_proc.dta} & Panel colapsado a nivel manzana + tratamiento & \texttt{3\_analysis\_did\_atributes.do} \\
\texttt{treat\_CEM.dta} & Muestra de 2018 con pesos CEM & \texttt{4\_analysis\_cem.do} \\
\texttt{treat\_weights.dta} & Pesos CEM para merge al panel & \texttt{4\_analysis\_cem.do} \\
\texttt{predios\_proc\_ai2.dta} & Pipeline completo sobre nueva AI & \texttt{6\_nueva\_ai.do} \\
\texttt{manzanas\_access.dta} & Índices de accesibilidad a nivel manzana & \texttt{7\_centralidades.do} \\
\bottomrule
\end{longtable}

% ============================================================
\section{Descripción de Outcomes}
% ============================================================

\subsection{Tablas de Regresión (formato \texttt{.docx}, generadas con \texttt{outreg2})}

\begin{table}[h!]
\centering
\begin{tabular}{@{}L{5cm} L{9.5cm}@{}}
\toprule
\textbf{Archivo} & \textbf{Contenido} \\
\midrule
\texttt{DID\_simple.docx} & ATT sobre avalúo catastral real — buffers 400m, 800m, 1200m \\
\texttt{DID\_heter\_dest.docx} & Heterogeneidad por destino económico \\
\texttt{DID\_heter\_estr.docx} & Heterogeneidad por estrato socioeconómico (residencial) \\
\texttt{DID\_heter\_dest\_res.docx} & ATT solo sobre predios residenciales \\
\texttt{DID\_heter\_tram.docx} & Heterogeneidad por tramo de la PLMB \\
\texttt{DID\_CEM.docx} & ATT con muestra emparejada por CEM \\
\texttt{DID\_atributes.docx} & Efecto sobre área construida (nivel manzana) \\
\texttt{DID\_atributes\_aim.docx} & Efecto sobre área construida — diseño AIM \\
\texttt{DID\_heter\_dest\_access.docx} & Heterogeneidad destino $\times$ accesibilidad \\
\texttt{DID\_heter\_estr\_access.docx} & Heterogeneidad estrato $\times$ accesibilidad (residencial) \\
\bottomrule
\end{tabular}
\caption{Tablas de regresión generadas por el pipeline.}
\end{table}

\newpage

\subsection{Gráficos de Coeficientes (formato \texttt{.pdf})}

\begin{table}[h!]
\centering
\begin{tabular}{@{}L{5cm} L{9.5cm}@{}}
\toprule
\textbf{Archivo} & \textbf{Contenido} \\
\midrule
\texttt{DiD\_windows\_1200.pdf} & DiD dinámico — avalúo catastral, buffer 1200m \\
\texttt{DiD\_windows\_com\_1200.pdf} & DiD dinámico — valor comercial, buffer 1200m \\
\texttt{DiD\_windows\_filter.pdf} & DiD dinámico — sin predios públicos \\
\texttt{DiD\_windows\_area.pdf} & DiD dinámico — área construida nivel manzana \\
\texttt{DiD\_windows\_tram.pdf} & DiD dinámico — por tramo (Tramo 1, 2, 3) \\
\texttt{DiD\_heter\_dest\_access.pdf} & ATT por destino económico dentro de cada quintil de accesibilidad \\
\bottomrule
\end{tabular}
\caption{Gráficos de efectos dinámicos generados por el pipeline.}
\end{table}

% ============================================================
\section{Requerimientos}
% ============================================================

Instalar los siguientes paquetes en Stata antes de ejecutar:

\begin{lstlisting}[language={}]
ssc install reghdfe
ssc install ftools
ssc install outreg2
ssc install cem
ssc install coefplot
\end{lstlisting}

% ============================================================
\section{Configuración de Rutas}
% ============================================================

Cada do-file define al inicio:

\begin{lstlisting}[language={}]
global dir_0 "C:\Users\USUARIO\\"
\end{lstlisting}

Los datos se esperan en:
\begin{lstlisting}[language={}]
${dir_0}OneDrive - Universidad de los andes\RA Andes - TIF\Datos\
\end{lstlisting}

Los outcomes se exportan a:
\begin{lstlisting}[language={}]
${dir_0}Documents\GitHub\TIF_PLMB\Datos\outcomes\
\end{lstlisting}

\begin{notebox}
Asegurarse de que solo una línea \texttt{global dir\_0} esté activa en cada archivo. La segunda línea (servidor de Uniandes) debe estar comentada con \texttt{//}.
\end{notebox}

% ============================================================
\section{Orden de Ejecución}
% ============================================================

Los do-files deben correrse en el siguiente orden. Los archivos 3a y 3b pueden ejecutarse en paralelo después del 2.

\begin{lstlisting}[language={}]
1_append_merge.do
    └── 2_analysis_did.do
            ├── 3_analysis_did_atributes.do
            ├── 3_analysis_did_heter.do
            ├── 4_analysis_cem.do
            └── 5_robustez.do
6_nueva_ai.do
    └── 7_centralidades.do
            └── 9_heter_access.do
8_robustez_control.do   (requiere predios_proc_ai2.dta de 6_nueva_ai.do)
\end{lstlisting}

\end{document}
